\begin{table}[H]
\centering
\caption{基于全局归一化理想点距离选取的代表性折中解}
\label{tab:representative-solutions}
\begin{tabular}{lccrrrr}
\toprule
算法 & 运行序号 & 随机种子 & AEP/(GWh/年) & LCOE & CVaR95/\% & 经验前沿贡献 \\
\midrule
CCGDE3 & 4 & 4 & 404.50 & 0.023157 & 6.62 & 22（17.74\%） \\
NSGA2  & 2 & 2 & 409.11 & 0.023112 & 4.61 & 102（82.26\%） \\
\bottomrule
\end{tabular}
\begin{minipage}{0.94\textwidth}
\footnotesize 注：三个目标先使用合并经验前沿的理想点和最差点归一化，再从各算法全部最终档案解中选择到理想点欧氏距离最小的解。
\end{minipage}
\end{table}

